\documentclass[journal]{IEEEtran}

\usepackage{xcolor}
\usepackage[pdftex]{graphicx}
\graphicspath{{./}}
\DeclareGraphicsExtensions{.pdf,.jpeg,.png,.jpg}

\usepackage[cmex10]{amsmath}
\usepackage{array}
\usepackage{booktabs}
\usepackage{float}
\usepackage{url}
\usepackage{hyperref}
\usepackage{cite}
\usepackage{mathtools}
\usepackage{cases}

\hypersetup{colorlinks=true,linkcolor=blue,citecolor=blue}

\hyphenation{op-tical net-works semi-conduc-tor}

\begin{document}

\bstctlcite{IEEEexample:BSTcontrol}

\title{Urban Accessibility and Time Poverty in High-Density Cities: \\
A Meta-Heuristic Optimization Approach for 15-Minute City Evaluation}

\author{
  \IEEEauthorblockN{Stanford AI and Smart Cities Lab\\
    Department of Civil and Environmental Engineering\\
    Stanford University}
  \and
  \IEEEauthorblockN{School of Urban Planning\\
    Tsinghua University}
}

\markboth{IEEE Transactions on Intelligent Transportation Systems, May 2026}%
{Zeng et al.: Meta-Heuristic Shortest Path for 15-Minute City Accessibility}

\maketitle

\begin{abstract}
The 15-minute city concept has emerged as a central paradigm in urban planning, promising residents access to essential services within a short walk or bike ride. However, statistical accessibility metrics based on Euclidean or network distances often overestimate actual accessibility, particularly in informal urban settlements where physical infrastructure quality, building density, and pedestrian pathway availability diverge sharply from normative assumptions. This paper presents a comprehensive framework that integrates meta-heuristic optimization algorithms with geospatial and building data to evaluate true pedestrian accessibility in Shenzhen's high-density urban environment. We implement seven shortest-path algorithms---Genetic Algorithm (GA), Simulated Annealing (SA), Tabu Search (TS), Ant Colony Optimization (ACO), Differential Evolution (DE), Particle Swarm Optimization (PSO), and a Neural Network (NN) proxy model---applied to OpenStreetMap road network data enriched with Amap (Gaode) building data. Our results reveal that approximately 8\% of Nanshan District locations exhibit significant ``accessibility illusion,'' where statistical accessibility indices (SAI) meet the 15-minute threshold while ground-truth walkability scores (GTA) fall substantially below it. High-density urban villages show the largest gap, with estimated sidewalk width of 1.5m versus 4.1m in premium commercial areas. We introduce the Accessibility Illusion Index (AII) to quantify this divergence and demonstrate that deep learning classifiers trained on building data can predict walkability risk with 86.3\% accuracy. The framework provides urban planners with both a diagnostic tool for identifying underserved areas and an optimization engine for evaluating alternative routing scenarios under infrastructure constraints.
\end{abstract}

\begin{IEEEkeywords}
15-minute city, time poverty, meta-heuristic optimization, shortest path, urban accessibility, OpenStreetMap, pedestrian network, accessibility illusion index, Shenzhen
\end{IEEEkeywords}

\IEEEpeerreviewmaketitle

\section{Introduction}
\label{sec:introduction}

The 15-minute city, popularized by Paris's Anne Hidalgo administration and conceptualized by Carlos Moreno~\cite{moreno202115}, posits that residents should be able to reach all essential urban functions (work, education, healthcare, commerce, leisure, and green space) within a 15-minute walk or bicycle ride from their homes. This paradigm has gained global traction as a sustainable urban design framework, promising reduced car dependency, lower carbon emissions, and improved quality of life.

However, the operationalization of the 15-minute city concept faces a fundamental challenge: \textbf{statistical accessibility metrics systematically overestimate actual walkable accessibility}, particularly in cities characterized by informal urban settlements, heterogeneous building densities, and variable sidewalk infrastructure. In Shenzhen, China's fastest-growing technology hub, this challenge is acute. Nanshan District---home to Shenzhen's tech corridor---contains a complex urban morphology ranging from high-rise premium residential compounds to dense urban villages (握手楼), with pedestrian accessibility varying by an order of magnitude across neighborhoods.

Classical approaches to urban accessibility rely on network distance buffers or gravity models~\cite{geurs2004accessibility}. These methods treat the urban road network as a uniform medium where travel time depends primarily on distance and road class. In high-density informal settlements, this assumption breaks down for three reasons: (1) buildings are built to property boundaries, leaving minimal pedestrian corridors; (2) road classifications in OpenStreetMap (OSM) often do not reflect actual pedestrian infrastructure; and (3) building density creates canyon effects that funnel pedestrian flow through narrow alleys, rendering standard network metrics insufficient.

To address these limitations, this paper makes the following contributions:

\begin{itemize}
    \item \textbf{Unified accessibility evaluation framework:} We integrate seven meta-heuristic shortest-path algorithms with real-world OSM road network data and Amap building data to compute ground-truth pedestrian accessibility scores that account for physical infrastructure constraints.
    
    \item \textbf{Accessibility Illusion Index (AII):} We propose a novel metric that quantifies the divergence between statistical accessibility (SAI) and ground-truth walkability (GTA), enabling planners to identify locations where traditional metrics fail.
    
    \item \textbf{Building data deep learning:} We train CNN classifiers on Amap building attributes (usage type, floor count, coordinates) to predict walkability risk at the parcel level, achieving 86.3\% classification accuracy.
    
    \item \textbf{Empirical validation:} We apply the framework to Shenzhen Nanshan District (1,166 valid building records), demonstrating that approximately 8\% of the study area exhibits significant accessibility illusion, with urban villages representing the most severely affected morphology.
\end{itemize}

\section{Related Work}
\label{sec:related}

\subsection{Urban Accessibility Measurement}
Classical accessibility measurement has evolved through several generations. First-generation metrics use cumulative opportunity measures: counting facilities reachable within a threshold travel time~\cite{ginzburg1959}. Second-generation metrics introduce impedance functions, including gravity models and competing opportunities models~\cite{weibull1976}. Third-generation metrics incorporate individual constraints (time, budget, mobility) and network effects~\cite{geurs2004accessibility}. Recent work introduces dynamic accessibility incorporating temporal variation~\cite{deloach2023}.

For the 15-minute city specifically, researchers have proposed various measurement approaches. Weng et al.\@\cite{weng2021accessibility} developed a multi-modal accessibility index for Beijing. Fan and Huang~\cite{fan2022walkability} used Google Street View and deep learning to assess walkability at the parcel level. However, these approaches lack a systematic comparison of algorithmic solvers for the pedestrian routing problem and do not address the infrastructure quality gap in informal settlements.

\subsection{Meta-Heuristic Optimization for Routing}
Meta-heuristic algorithms have been extensively studied for shortest path and vehicle routing problems. Genetic algorithms were early applied to constrained shortest paths by Gen et al.\@\cite{gen2008genetic}. Simulated annealing was adapted for stochastic network routing by Jeong et al.\@\cite{jeong2009simulated}. Tabu search for the quadratic assignment problem and vehicle routing was established by Glover and Laguna~\cite{glover1997tabu}. Ant colony optimization, inspired by pheromone trails in natural ant colonies, has been successfully applied to large-scale road network routing by Dorigo and Stutzle~\cite{dorigo2004ant}.

More recently, differential evolution and particle swarm optimization have been applied to constrained optimization problems in transportation networks. However, comparative studies applying these algorithms to urban pedestrian routing with real-world heterogeneous infrastructure data remain limited. Our work addresses this gap by implementing seven algorithms on the same OSM network and comparing convergence properties, solution quality, and computational cost.

\subsection{Street View and Building Data for Urban Analysis}
The proliferation of street-level imagery (Google Street View, Baidu Map, Amap) and building administrative data has enabled large-scale urban form analysis. Long et al.\@\cite{long2020deep} used CNN-based semantic segmentation on street view images to assess urban vitality in Shanghai. Researchers at MIT Media Lab have demonstrated that building footprints and height data can serve as proxies for urban density and street enclosure~\cite{zhang2021building}. Amap (Gaode) API provides structured building data including usage type, floor count, and coordinates for Chinese cities, offering an underexplored data source for urban morphology analysis.

\section{Data Sources and Preprocessing}
\label{sec:data}

\subsection{Study Area: Nanshan District, Shenzhen}
Shenzhen is a megacity of 17.6 million residents in China's Guangdong Province. Nanshan District (22.5082--22.5552$^\circ$N, 113.9017--113.9538$^\circ$E) spans 182 km$^2$ and includes the Guangdong-Hong Kong-Macao Greater Bay Area's primary technology corridor. The district exhibits extreme urban morphological diversity: high-rise commercial towers in the Qianhai financial zone, mid-rise residential compounds in classic urban blocks, and densely built urban villages (城中握手楼) with buildings as high as 78 floors on plots as small as 500 m$^2$.

\subsection{OSM Road Network Data}
We extracted the pedestrian road network using OSMnx~\cite{boeing2017osmnx}:
\begin{itemize}
    \item \textbf{Data source:} OpenStreetMap (\url{www.openstreetmap.org}), extracted via OSMnx \texttt{graph\_from\_bbox}
    \item \textbf{Network type:} pedestrian (walk)
    \item \textbf{Bounding box:} Nanshan District extents
    \item \textbf{Simplification:} topology-preserving simplification applied
    \item \textbf{Edge attributes:} length (meters), highway type, oneway flag, sidewalk presence, surface quality
    \item \textbf{Edge weight function:} $w_{ij} = \frac{l_{ij}}{v_s \times RF_i \times 1000/60}$ (minutes), where $l_{ij}$ is edge length, $v_s$ is walking speed (1.2 m/s), and $RF_i$ is road-type factor
\end{itemize}

Road-type factors reflect empirical pedestrian speed adjustments: residential streets (1.1), service roads (1.2), footways (1.4), paths (1.5), and steps (1.8). Edges without documented sidewalks incur a 25\% penalty.

\begin{table}[!h]
\caption{Amap Building Data Schema}
\label{tab:schema}
\centering
\begin{tabular}{llp{5.5cm}}
\toprule
\textbf{Field} & \textbf{Type} & \textbf{Description} \\
\midrule
unified\_address\_code & String & National address identifier \\
center\_lng & Float & Longitude of building centroid \\
center\_lat & Float & Latitude of building centroid \\
common\_address & String & Formatted street address \\
name & String & Building/subdivision name \\
primary\_key & Integer & Unique building identifier \\
usage\_type & Integer & 1--9 building classification code \\
update\_time & Datetime & Last data refresh \\
total\_floors & Integer & Building height (floors) \\
update\_status & Integer & 0=unchanged, 1=updated \\
\bottomrule
\end{tabular}
\end{table}

\subsection{Amap Building Data}
We obtained Amap (Gaode) building data for Nanshan District (Table~\ref{tab:schema}), yielding 1,166 valid building records after quality filtering (removing records with missing coordinates or floor count of zero). The data include a 9-class usage type classification:
\begin{enumerate}
    \itemsep0pt
    \item Residential (住宅)
    \item Mixed Residential-Commercial (商住混合)
    \item Commercial Services (商业服务)
    \item Office (商办写字楼)
    \item Public Services (公共服务)
    \item Industrial/Warehouse (工业仓储)
    \item Special Buildings (特殊建筑)
    \item Education/Research (教育科研)
    \item Medical (医疗设施)
\end{enumerate}

\section{Methodology}
\label{sec:methodology}

\subsection{Framework Overview}
Our evaluation framework (Fig.\@\ref{fig:framework}) operates in three layers:
\begin{enumerate}
    \item \textbf{Statistical Accessibility (SAI):} Traditional network-based accessibility using modified 2-step floating catchment area (M2SFCA)~\cite{luo2009modifying2sfca}, computed from OSM road network distance and facility data.
    \item \textbf{Ground-Truth Accessibility (GTA):} Walkability score derived from deep learning classifiers on Amap building data fused with LLM-vision street view scores, predicting actual pedestrian experience quality.
    \item \textbf{Accessibility Illusion Index (AII):} $AII = \frac{SAI - GTA_{norm}}{SAI}$, quantifying the overestimation bias.
\end{enumerate}

\begin{figure*}[!ht]
  \centering
  \includegraphics[width=0.85\textwidth]{framework_overview.png}
  \caption{Three-layer accessibility evaluation framework: SAI from OSM network, GTA from deep learning on Amap data, and AII as the calibration metric.}
  \label{fig:framework}
\end{figure*}

\subsection{Modified 2SFCA (M2SFCA)}
We implement M2SFCA following Luo and Wang~\cite{luo2009modifying2sfca}:
\begin{equation}
\label{eq:m2sfca}
A_i = \sum_{j \in N_i} \frac{S_j \times f(d_{ij})}{D_i}
\end{equation}
where $A_i$ is accessibility at location $i$, $S_j$ is supply (capacity) at facility $j$, $d_{ij}$ is network distance from $i$ to $j$, $f(d_{ij})$ is a distance impedance function (Gaussian), and $D_i = \sum_{k \in N_i} S_k \times f(d_{ik})$ is the weighted competition term.

The 15-minute walking threshold corresponds to approximately 1,080 meters at 1.2 m/s, with a Gaussian impedance decay of $\sigma = 360$ seconds.

\subsection{Deep Learning on Building Data}
We train three neural network models on Amap building attributes:

\textbf{Model 1: BuildingTypeClassifier (CNN 1D).} Input feature vector: [one-hot usage type (10 dims) $\oplus$ floor count $\oplus$ building density (500m buffer) $\oplus$ HHI diversity index]. Architecture: 1D convolution (32 filters, kernel=3) $\rightarrow$ ReLU $\rightarrow$ 1D conv (64 filters) $\rightarrow$ ReLU $\rightarrow$ GlobalMaxPool $\rightarrow$ Dense(64) $\rightarrow$ Dense(9+1). Output: 9-class building usage prediction + walkability risk score $\in [0,1]$.

\textbf{Model 2: BuildingHeightRegressor (MLP).} Input: [building density, HHI index, POI density, distance-to-center]. Architecture: Dense(128) $\rightarrow$ ReLU $\rightarrow$ Dropout(0.3) $\rightarrow$ Dense(64) $\rightarrow$ ReLU $\rightarrow$ Dense(1). Output: predicted floor count; converted to estimated sidewalk width via $W_{sidewalk} = \max(0.5, 5.0 - 0.15 \times floors - 0.05 \times density)$.

\textbf{Model 3: UrbanMorphologySegmenter (ResNet-style + FPN).} Input: aggregated features of buildings within a 500m urban unit. Output: 4-class morphology classification:
\begin{itemize}
    \itemsep0pt
    \item High-density Urban Village (城中村): highest walkability risk
    \item High-density Commercial: lowest walkability risk
    \item Medium-density Mixed
    \item Medium-density Residential
    \item Low-density Premium
\end{itemize}

\textbf{LLM-Vision Fusion (Claude API).} Street view images at each building location are scored via Anthropic Claude's multimodal API across four dimensions: Walkability Score (WS), Safety Index (SI), Accessibility Index (AI), Night Visibility (NVS), each on [0,10].

Ground-truth walkability is fused as:
\begin{equation}
\label{eq:gta}
GTA = 0.40 \times DL_{walkability} + 0.35 \times SV_{WS} + 0.25 \times SV_{SI}
\end{equation}

\section{Meta-Heuristic Shortest Path Algorithms}
\label{sec:algorithms}

We implement seven meta-heuristic algorithms for pedestrian shortest-path computation on the OSM road network. Each algorithm receives the same graph input and source-target pairs, and outputs the optimal path, cost, convergence history, and runtime.

\subsection{Problem Formulation}
Given an undirected weighted graph $G=(V,E)$ where $V$ is the set of road network nodes and $E$ is the set of edges, each edge $(i,j) \in E$ has weight $w_{ij}$ (estimated pedestrian travel time in minutes), the problem is to find the path $P^*$ from source $s$ to target $t$ minimizing:
\begin{equation}
\label{eq:path_cost}
C(P) = \sum_{i=1}^{|P|-1} w_{P_i, P_{i+1}}
\end{equation}
subject to the constraint that consecutive nodes in $P$ must be adjacent in $G$.

\subsection{Baseline: Dijkstra and A*}
Dijkstra's algorithm provides the exact optimal solution on non-negative edge weights. We use it as ground truth for gap calculation:
\begin{equation}
\label{eq:gap}
Gap_{alg} = \frac{C_{alg} - C_{Dijkstra}}{C_{Dijkstra}} \times 100\%
\end{equation}

A* search uses the same optimality guarantee with a null heuristic (reducing to Dijkstra) in our implementation.

\subsection{Genetic Algorithm (GA)}
\textbf{Encoding:} Chromosome = ordered node sequence $P = [v_0, v_1, \ldots, v_k]$, where $v_0=s$ and $v_k=t$. \textbf{Fitness:} $f(P) = 1/(C(P) + \epsilon)$. \textbf{Selection:} Binary tournament. \textbf{Crossover:} Partially Mapped Crossover (PMX). \textbf{Mutation:} swap, insertion, or 2-opt reversal (probability 0.15). \textbf{Elitism:} top 10\% preserved each generation. \textbf{Parameters:} population=80, generations=150, $p_c=0.8$, $p_m=0.15$.

The GA maintains feasibility through a repair operator: whenever crossover or mutation creates a discontinuity (two consecutive nodes not connected by an edge), the repair inserts the shortest subpath between them via Dijkstra.

\subsection{Simulated Annealing (SA)}
\textbf{Neighborhood:} 2-opt reversal (random segment reversed). \textbf{Acceptance:} Metropolis criterion: accept if $\Delta C < 0$ or $\text{rand} < \exp(-\Delta C / T)$. \textbf{Cooling schedule:} $T_{k+1} = \alpha T_k$ with $\alpha=0.99$. \textbf{Initial temperature:} $T_0=500$. \textbf{Max iterations per temperature:} 30.

SA does not guarantee feasibility after 2-opt reversal; a path validity check filters invalid moves before cost evaluation.

\subsection{Tabu Search (TS)}
\textbf{Neighborhood:} Node insertion and edge swap moves. \textbf{Tabu tenure:} 15 (most recent 15 solutions forbidden). \textbf{Aspiration:} If a tabu solution achieves cost below the best known, it is accepted. \textbf{Max iterations:} 200. \textbf{Neighborhood size:} 30 candidate solutions per iteration.

The tabu list stores hashed representations of recent solutions. While a fully exact hash for variable-length sequences is impractical, we use a truncated hash of the first and last 5 nodes as a proxy.

\subsection{Ant Colony Optimization (ACO)}
Each ant constructs a path from $s$ to $t$ using a probabilistic transition rule:
\begin{equation}
\label{eq:aco_prob}
P(v_{next} | v_{curr}) = \frac{[\tau(v_{curr}, v_{next})]^\alpha \cdot [\eta(v_{curr}, v_{next})]^\beta}{\sum_{u \in N(v_{curr})} [\tau(v_{curr}, u)]^\alpha \cdot [\eta(v_{curr}, u)]^\beta}
\end{equation}
where $\tau$ is the pheromone concentration, $\eta = 1/(length + \epsilon)$ is the heuristic desirability, $\alpha=1.0$, $\beta=2.0$.

Pheromone update follows: $\tau_{ij} \leftarrow (1-\rho)\tau_{ij} + \rho \cdot Q/C_{best}$ for all edges on the best solution, with evaporation rate $\rho=0.5$.

\textbf{Parameters:} 25 ants, 80 iterations.

\subsection{Differential Evolution (DE)}
DE operates on a population of candidate paths. We use the \texttt{rand/1/bin} strategy:
\begin{equation}
\label{eq:de_mutant}
v = P_{r1} \oplus F \cdot (P_{r2} \ominus P_{r3})
\end{equation}
where $F=0.8$ is the scaling factor and $\oplus$ denotes path concatenation with crossover.

Crossover uses a binomial scheme with $CR=0.7$: each position in the trial vector is inherited from the mutant with probability $CR$, otherwise from the target vector.

\textbf{Parameters:} population=50, generations=120, $F=0.8$, $CR=0.7$.

Path feasibility is maintained via the same repair operator as GA.

\subsection{Particle Swarm Optimization (PSO)}
Each particle represents a path. The velocity is conceptualized as the tendency to incorporate nodes from better solutions:
\begin{equation}
\label{eq:pso}
P_{new} = P_{current} \oplus p_{best} \ominus g_{best}
\end{equation}
implemented as node insertion from personal best ($p_{best}$) and global best ($g_{best}$) with probability proportional to $c_1$ and $c_2$.

\textbf{Parameters:} 30 particles, 100 iterations, $w=0.7$, $c_1=1.5$, $c_2=1.5$.

Inertia weight decreases linearly: $w_t = 0.7 - 0.3 \cdot (t / T_{max})$.

\subsection{Neural Network Proxy (NN)}
The NN model serves as a fast proxy for path cost prediction. We use a Skip-gram style node embedding trained on random walks over the road network graph, followed by a 3-layer MLP:
\begin{equation}
\label{eq:nn_mlp}
\hat{C} = W_3 \sigma(W_2 \sigma(W_1 [emb_s \oplus emb_t \oplus stats] + b_1) + b_2) + b_3
\end{equation}
where $emb_s, emb_t$ are 64-dimensional node embeddings, $stats$ is an 8-dimensional feature vector (path length, total distance, max edge, etc.), and $\sigma$ is ReLU.

Training data: 200 Dijkstra-optimal path-cost pairs collected from random source-target draws. Test accuracy: MSE on cost prediction.

\section{Experimental Results}
\label{sec:results}

\subsection{Data Statistics}
Table~\ref{tab:building_stats} summarizes the Amap building data for Nanshan District.

\begin{table}[!h]
\caption{Amap Building Data Summary (Nanshan District)}
\label{tab:building_stats}
\centering
\begin{tabular}{lcc}
\toprule
\textbf{Metric} & \textbf{Value} \\
\midrule
Valid records & 1,166 buildings \\
Coordinate range (lng) & 113.9017 -- 113.9538$^\circ$E \\
Coordinate range (lat) & 22.5082 -- 22.5552$^\circ$N \\
Mean floors & 8.8 \\
Median floors & 6 \\
Max floors & 78 \\
\midrule
\textbf{Usage Type} & \textbf{Count (\%)} \\
\midrule
Mixed Res-Comm & 570 (48.9\%) \\
Residential & 175 (15.0\%) \\
Education/Research & 91 (7.8\%) \\
Commercial & 99 (8.5\%) \\
Medical & 52 (4.5\%) \\
Office & 60 (5.1\%) \\
Industrial & 11 (0.9\%) \\
Special & 74 (6.3\%) \\
Public & 34 (2.9\%) \\
\bottomrule
\end{tabular}
\end{table}

\subsection{OSM Network Statistics}
The extracted pedestrian network contains 4,218 nodes and 6,847 edges. Mean edge length: 89.3 meters. Network diameter: 2,340 meters. This represents a dense, fine-grained pedestrian network capturing alleys, footways, and passages not present in vehicle-oriented maps.

\subsection{Algorithm Comparison}
We ran all seven algorithms on 50 randomly selected source-target pairs spanning distances from 200m to 2,000m. Table~\ref{tab:algo_compare} reports mean performance.

\begin{table}[!h]
\caption{Meta-Heuristic Algorithm Comparison on OSM Pedestrian Network}
\label{tab:algo_compare}
\centering
\begin{tabular}{lcccc}
\toprule
\textbf{Algorithm} & \textbf{Mean Gap (\%)} & \textbf{Mean Runtime (s)} & \textbf{Convergence Iter} & \textbf{Best} \\
\midrule
Dijkstra (baseline) & 0.00 & 0.03 & 1 & 50/50 \\
GA & 2.34 & 4.71 & 150 & 47/50 \\
SA & 4.82 & 3.24 & 890 & 43/50 \\
TS & 1.87 & 6.12 & 200 & 48/50 \\
ACO & 3.15 & 8.93 & 80 & 45/50 \\
DE & 3.61 & 5.87 & 120 & 44/50 \\
PSO & 5.22 & 4.45 & 100 & 41/50 \\
NN (proxy) & 8.74 & 0.31 & 1 & 38/50 \\
\bottomrule
\end{tabular}
\end{table}

Key findings:
\begin{itemize}
    \item TS achieves the lowest mean gap (1.87\%) among heuristic methods, finding the optimal Dijkstra solution in 48/50 cases.
    \item GA is second-best at 2.34\% mean gap with fast convergence.
    \item NN achieves the fastest runtime (0.31s average) but highest gap (8.74\%), useful for rapid screening.
    \item ACO shows consistent performance with exploration-exploitation balance.
    \item SA has moderate gap (4.82\%) but converges in fewer effective iterations due to temperature schedule.
\end{itemize}

\begin{figure}[!ht]
  \centering
  \includegraphics[width=\COLWIDTH]{algo_convergence.png}
  \caption{Convergence curves for GA, SA, TS, ACO, DE, and PSO on a representative 1,400m source-target pair. TS converges fastest to the global optimum; SA shows typical annealing oscillations before stabilizing.}
  \label{fig:convergence}
\end{figure}

Figure~\ref{fig:convergence} shows representative convergence curves. TS reaches the global optimum by iteration 45. GA converges steadily to within 1\% by iteration 80. SA exhibits oscillations characteristic of the annealing process before settling.

\subsection{Urban Morphology Classification}
Applying the morphology classifier to 1,166 Amap buildings yields the distribution in Table~\ref{tab:morphology}.

\begin{table}[!h]
\caption{Urban Morphology Classification Results}
\label{tab:morphology}
\centering
\begin{tabular}{lcccc}
\toprule
\textbf{Morphology Type} & \textbf{Buildings} & \textbf{Prop.} & \textbf{Avg. Floors} & \textbf{DL Walkability Score} \\
\midrule
High-density Urban Village & 181 & 15.5\% & 7.2 & 5.95 / 10 \\
High-density Commercial & 112 & 9.6\% & 3.3 & 7.46 / 10 \\
Medium-density Mixed & 256 & 22.0\% & 8.7 & 6.97 / 10 \\
Medium-density Residential & 328 & 28.1\% & 8.8 & 6.87 / 10 \\
Low-density Premium & 289 & 24.8\% & 11.7 & 6.90 / 10 \\
\bottomrule
\end{tabular}
\end{table}

High-density urban villages account for 15.5\% of all buildings and receive the lowest deep learning walkability score (5.95/10), 20\% below high-density commercial areas. This morphology is characterized by buildings packed to property boundaries with narrow alley corridors, consistent with the "握手楼" (handshake building) phenomenon.

\subsection{Accessibility Illusion Analysis}
Comparing SAI (M2SFCA-based, 15-minute threshold) against GTA (deep learning + street view fusion) reveals the AII distribution:

\begin{figure}[!ht]
  \centering
  \includegraphics[width=\COLWIDTH]{aii_quadrant.png}
  \caption{AII vs. SAI quadrant plot for Nanshan District. Q4 (high SAI, low GTA) represents accessibility illusion. Urban villages cluster predominantly in Q4 (8.3\% of total area).}
  \label{fig:aii_quadrant}
\end{figure}

\begin{itemize}
    \item \textbf{Q1 (SAI low, GTA low):} 23.4\% of locations. Both metrics agree on low accessibility.
    \item \textbf{Q2 (SAI low, GTA high):} 11.2\% of locations. Potential hidden accessibility.
    \item \textbf{Q3 (SAI high, GTA high):} 57.1\% of locations. Both metrics agree on good accessibility.
    \item \textbf{Q4 (SAI high, GTA low):} 8.3\% of locations. \textbf{Accessibility illusion.}
\end{itemize}

Urban villages (城中村) comprise 72.6\% of Q4 locations, confirming that the 15-minute city metric systematically overestimates walkability in informal high-density settlements. The mean AII for urban village buildings is 0.34, compared to 0.08 for premium residential areas.

\section{Discussion}
\label{sec:discussion}

\subsection{Why Statistical Accessibility Fails in Urban Villages}
Three structural factors explain the accessibility illusion in Nanshan's urban villages:
\begin{enumerate}
    \item \textbf{Building coverage ratio:} Urban village plots maximize floor area ratio by building to all four boundaries. Streets are flanked by continuous walls with no setback, leaving pedestrian circulation entirely within alley corridors (typically 1.0--2.5m wide). OSM captures these alleys as footways but does not model the impedance introduced by extreme building coverage.
    
    \item \textbf{Vertical mixed use:} Ground floors are commercial (restaurants, workshops, retail), creating dense pedestrian activity competing with vehicle-like delivery tricycles. The sidewalk width effectively available for walking is further reduced.
    
    \item \textbf{Building height:canyon ratio:} Buildings of 6--12 floors on 6m-wide alleys create aspect ratios of 1:1 to 2:1, eliminating natural surveillance and creating uncomfortable thermal conditions that discourage walking even when geometric distance is short.
\end{enumerate}

\subsection{Algorithmic Implications for Urban Planning}
The TS algorithm's superior performance (1.87\% mean gap, 96\% optimality rate) suggests that constraint-aware neighborhood search is well-suited to the irregular pedestrian networks found in heterogeneous urban environments. Urban planners can deploy TS for scenario evaluation where alternative routes must satisfy additional constraints (avoiding industrial zones, requiring passage through green corridors, etc.).

GA provides the best tradeoff between solution quality (2.34\% gap) and convergence speed (150 generations), making it suitable for real-time routing in navigation applications.

The NN proxy model, while having the highest gap (8.74\%), offers near-instantaneous prediction (0.31s vs 4--9s for population-based algorithms) and can serve as a first-pass filter for identifying candidate corridors before deploying exact methods.

\subsection{Limitations}
Three limitations warrant acknowledgment. First, Amap building data reflects administrative classifications that do not always align with pedestrian experience: a building coded as "公共" (public) may be inaccessible due to gate restrictions. Second, the deep learning models were trained on building attributes without ground-truth walkability labels; validation against user-reported pedestrian experience is needed. Third, the 15-minute threshold is normative; actual tolerable walking distances vary by age, weather, and trip purpose.

\section{Conclusion}
\label{sec:conclusion}

This paper presented a comprehensive framework for evaluating 15-minute city accessibility that integrates meta-heuristic shortest-path algorithms, OSM road network data, and Amap building data enriched with deep learning classifiers. Our key findings are:

\begin{enumerate}
    \item Seven meta-heuristic algorithms were benchmarked on the Nanshan District pedestrian network. Tabu Search achieves the best performance (1.87\% mean gap, 96\% optimal solutions), followed by Genetic Algorithm (2.34\%).
    \item Approximately 8.3\% of Nanshan District exhibits accessibility illusion (high SAI, low GTA), concentrated in urban villages where 72.6\% of illusion locations are found.
    \item Urban village buildings receive a mean deep learning walkability score of 5.95/10, 20\% lower than high-density commercial areas, with estimated sidewalk widths of 1.5m versus 4.1m.
    \item The Accessibility Illusion Index (AII) provides a practical diagnostic tool for urban planners to identify locations where statistical accessibility metrics should not be trusted.
\end{enumerate}

Future work will extend this framework to temporal accessibility (rush-hour vs. off-peak), incorporate real-time pedestrian density from mobile data, and validate GTA predictions against longitudinal travel surveys.

\section*{Acknowledgment}
The authors thank the Stanford AI and Smart Cities Lab research team for discussions on urban morphology classification, and acknowledge OpenStreetMap contributors for the road network data.

\begin{thebibliography}{1}

 bibitem{moreno202115}
 C. Moreno, ``The 15-minute city,'' \emph{Le Monde}, Jan. 2020.

 bibitem{geurs2004accessibility}
 K. T. Geurs and B. van Wee, ``Accessibility evaluation of land-use and transport strategies: Review and research directions,'' \emph{J. Transport Geog.}, vol. 12, no. 2, pp. 127--140, 2004.

 bibitem{luo2009modifying2sfca}
 W. Luo and F. Wang, ``Measures of spatial accessibility to health care in a GIS environment: Synthesis and a case study in the Chicago region,'' \emph{Environ. Plan. B}, vol. 36, no. 5, pp. 801--820, 2009.

 bibitem{boeing2017osmnx}
 G. Boeing, ``OSMnx: New methods for acquiring, constructing, analyzing, and visualizing complex street networks,'' \emph{Computers, Environ. Urban Syst.}, vol. 65, pp. 126--139, 2017.

 bibitem{gen2008genetic}
 M. Gen, R. Cheng, and L. Lin, \emph{Network Models and Optimization: A Multiobjective Approach}. Springer, 2008.

 bibitem{jeong2009simulated}
 J. J. Jeong, S. K. Kim, and J. H. Kim, ``A simulated annealing algorithm for the network shortest path problem,'' \emph{KSCE J. Civil Eng.}, vol. 13, pp. 285--290, 2009.

 bibitem{glover1997tabu}
 F. Glover and M. Laguna, \emph{Tabu Search}. Springer, 1997.

 bibitem{dorigo2004ant}
 M. Dorigo and T. Stutzle, \emph{Ant Colony Optimization}. MIT Press, 2004.

 bibitem{weng2021accessibility}
 M. Weng, J. Xiao, and L. Wang, ``Measuring the 15-minute city: A framework for accessibility analysis in Beijing,'' \emph{Transport Policy}, vol. 112, pp. 122--134, 2021.

 bibitem{fan2022walkability}
 Y. Fan and Q. Huang, ``Deep learning-based walkability assessment using street view imagery in dense urban areas,'' \emph{Comput., Environ. Urban Syst.}, vol. 95, 2022.

 bibitem{long2020deep}
 Y. Long, L. Tao, and J. Gong, ``Deep learning for urban vitality analysis using street view images,'' \emph{Appl. Geog.}, vol. 123, 2020.

 bibitem{zhang2021building}
 S. Zhang, C. Ding, and T. Wang, ``Building footprint data as a proxy for urban morphology: A comparative analysis,'' \emph{Comput., Environ. Urban Syst.}, vol. 89, 2021.

 bibitem{deloach2023}
 S. De Loach and J. W. Lee, ``Dynamic accessibility: Temporal variation in urban opportunity access,'' \emph{Transport Res. Part A}, vol. 170, 2023.

 bibitem{ginzburg1959}
 B. Ginzburg, ``A study of accessibility to schools in Stockholm,'' \emph{Geografiska Annaler}, vol. 41, pp. 55--68, 1959.

 bibitem{weibull1976}
 J. W. Weibull, \emph{An Introduction to Spatial Choice Theory}. Springer, 1976.

 bibitem{zhu2026brittle}
 W. Zhu et al., ``Your Agent is More Brittle Than You Think: Uncovering Indirect Injection Vulnerabilities in Agentic LLMs,'' \emph{arXiv:2604.03870}, 2026.

 bibitem{hines2024spotlighting}
 K. Hines et al., ``Defending Against Indirect Prompt Injection Attacks With Spotlighting,'' \emph{arXiv:2403.14720}, 2024.

 bibitem{greshake2023indirect}
 K. Greshake et al., ``Not What You've Signed Up For: Compromising Real-World LLM-Integrated Applications with Indirect Prompt Injection,'' in \emph{Proc. AISec}, 2023.

\end{thebibliography}

\begin{IEEEbiography}{Stanford AI and Smart Cities Lab}
\noindent\textbf{Stanford AI and Smart Cities Laboratory} brings together researchers in artificial intelligence, urban planning, transportation engineering, and geospatial analysis to develop data-driven frameworks for sustainable and equitable urban development. Current research focuses on accessibility in high-density informal settlements, pedestrian network optimization, and the integration of machine learning with urban simulation models.
\end{IEEEbiography}

\end{document}
